% ~6 min
\partframe{Part III}{Program Calculation and Extraction}

\begin{frame*}
  \frametitle{Calculating a Program with \texttt{Derive}}
  \rocqblock{qsort-derive.v}
  \ocamlblock{qsort-extr.ml}
  \speakernote{Derive leaves qsort as an unknown; the rewrites build it,
    and finished fills it in, reduced to the code to extract. The last
    rewrite switches to a continuation-passing coalgebra, so the code has
    no closures. qsort is transparent, so it extracts; qsort\_calc is the
    proof that it meets the specification. The OCaml is what one would
    write by hand.}
\end{frame*}

\begin{frame*}
  \frametitle{Fusing a Map into Quicksort}
  \rocqblock{qsort2-derive.v}
  \ocamlblock[before upper = {\fvset{fontsize = \footnotesize}}]{qsort2.ml}
  \speakernote{The opening example. One rewrite chain: fuse the map into
    merge (left fusion), then deforest. The extracted code doubles each
    pivot in a single traversal.}
\end{frame*}

\begin{frame*}[vertical-alignment = top]
  \frametitle{Quicksort, Step by Step}
  \vspace*{4mm}
  \begin{onlyenv}<1>
    \fvset{lastline = 1}%
    \rocqblock{qsort-steps.v}
  \end{onlyenv}
  \begin{onlyenv}<2>
    \fvset{lastline = 3}%
    \rocqblock{qsort-steps.v}
  \end{onlyenv}
  \begin{onlyenv}<3>
    \fvset{lastline = 5}%
    \rocqblock{qsort-steps.v}
  \end{onlyenv}
  \begin{onlyenv}<4->
    \fvset{lastline = 7}%
    \rocqblock{qsort-steps.v}
  \end{onlyenv}
  \par
  \speakernote{Each step: one law, the term Rocq shows after it \\
    map\_merge\_fusion: doubling moves inside merge \\
    hylo\_ana, deforest: the tree between unfold and fold disappears (one step) \\
    hylo\_cps\_hylo: continuation-passing coalgebra, so the code has no closures}
\end{frame*}

\begin{frame*}
  \frametitle{Coalgebras Need Not Follow the Data}
  \rocqblock{digits.v}
  \begin{center}
    \begin{tikzpicture}[alt={The digits coalgebra on 13 in base 2: 13
        gives digit 1 and seed 6, 6 gives digit 0 and seed 3, 3 gives
        digit 1 and seed 1, 1 gives digit 1 and seed 0, and 0 stops.},
      node distance = 12mm]
      \node[box] (a) {$13$};
      \node[box, right = of a] (b) {$6$};
      \node[box, right = of b] (c) {$3$};
      \node[box, right = of c] (d) {$1$};
      \node[box, right = of d] (e) {$0$};
      \draw[->] (a) -- node[above, note] {1} (b);
      \draw[->] (b) -- node[above, note] {0} (c);
      \draw[->] (c) -- node[above, note] {1} (d);
      \draw[->] (d) -- node[above, note] {1} (e);
      \node[note, right = 3mm of e] {digits of $13 = 1101_2$};
    \end{tikzpicture}
  \end{center}
  \ocamlblock{fastpow.ml}
  \speakernote{digits splits a number, no list is ever built. Power as a
    fold over the binary digits; deforestation gives fast exponentiation.
    The same laws can also move the map to the other side and introduce a
    list (sum of squares example in the library).}
\end{frame*}

\begin{frame*}
  \frametitle{Challenge: Inlining for Extraction}
  Extracting the unreduced \code{hylo merge tsplit}:
  \ocamlblock{raw-extr.ml}
  \callout<2->{14.5,8.2}{dictionaries at run time}{dict}
  \callout<3->{14,2.6}{a record per call}{rec}
  \speakernote{Shown partially. Type-safe, but generic: setoid and
    container dictionaries are passed around at run time, and a record is
    built and taken apart at every call. The fix is now in the library:
    finished reduces the derived program (unfolding the unfold\_code hints)
    and CPS coalgebras remove the closures, so a plain Extraction of qsort
    gives the code from before.}
\end{frame*}

\begin{frame}
  \frametitle{Since the Paper: Usability}
  \begin{itemize}
    \item<1-> Properness and termination: \textbf{instance search}.
    \item<2-> Rewriting modulo associativity: \code{arewrite}.
    \item<3-> Diverging instance search: \code{\#[mode = "!"]}.
  \end{itemize}
  \speakernote{Before: every function needed a properness proof, every
    coalgebra a termination proof; now found by instance search, and a
    tactic proves properness out of container values. Calculations are
    long chains of compositions; rewrite needs a syntactic match up to
    associativity; arewrite finds the law anywhere in the chain and proves
    the step by conversion. An unannotated statement made Rocq search for
    setoid ?A; some instances recurse on a fresh evar, so the search never
    ended and ate all memory; the mode annotation stops it.}
\end{frame}
